For $0\le h\le1$, let $d=1-h$ be the distance of the chord above the centre and put $\theta=\arccos d$ in radians. The sector angle is $2\theta$, hence its area is $\theta$; the triangle between the two radii has area $d\sqrt{1-d^2}$. Subtracting gives $A(h)=\arccos(1-h)-(1-h)\sqrt{2h-h^2}$. For $1<h\le2$, the complementary lower cap has height $2-h$, so $A(h)=\pi-A(2-h)$, which simplifies to the same formula. Thus $\boxed{A(h)=\arccos(1-h)-(1-h)\sqrt{2h-h^2}}$ for the entire interval, including endpoint areas $0$ and $\pi$.
\begin{asy}
texpreamble("\usepackage{amsmath,amssymb}");
size(240);
defaultpen(fontsize(9pt));
pair O=(0,0),L=(-0.8660254037844386,0.5),R=(0.8660254037844386,0.5);filldraw(arc(O,1,30.000000000000004,150.0)--cycle,lightblue);draw(circle(O,1));draw(O--L,dashed);draw(O--R,dashed);draw((1.15,0.5)--(1.15,1),Arrows);label("$h$",(1.15,0.75),E);dot(O);label("O",O,S);label("L",L,W);label("R",R,E);
\end{asy}
\begin{asy}
texpreamble("\usepackage{amsmath,amssymb}");
size(240);
defaultpen(fontsize(9pt));
pair O=(0,0),L=(-0.8660254037844386,-0.5),R=(0.8660254037844386,-0.5);filldraw(arc(O,1,-30.000000000000004,210.0)--cycle,lightblue);draw(circle(O,1));draw(O--L,dashed);draw(O--R,dashed);draw((1.15,-0.5)--(1.15,1),Arrows);label("$h$",(1.15,0.25),E);dot(O);label("O",O,S);label("L",L,W);label("R",R,E);
\end{asy}
